\documentclass[11pt,answers]{exam}

\usepackage{etex}
\usepackage{amssymb,amsmath,multicol} %<-- InWorksheetExam1 i also have fancyhdr,

\usepackage[metapost]{mfpic}
\usepackage[pdftex]{graphicx}
\usepackage{tabu}

\usepackage{pst-plot}
\usepackage{pgfplots}
\pgfplotsset{compat=1.9}

\usepackage{tikz}
\usepackage{tkz-2d}
\usepackage{tkz-base}
\usetikzlibrary{calc}
\usetikzlibrary{arrows}

\usepackage{systeme}

\usepackage[inline]{enumitem}
\usepackage{refcount}%<-- non in WorksheetExam1

\usepackage{pstricks-add,pst-eucl}
\usepackage{systeme}
\usepackage{setspace}
\usepackage{multicol}


\usepackage[inline]{enumitem}   
\makeatletter
% This command ignores the optional argument for itemize and enumerate lists
\newcommand{\inlineitem}[1][]{%
\ifnum\enit@type=\tw@
    {\descriptionlabel{#1}}
  \hspace{\labelsep}%
\else
  \ifnum\enit@type=\z@
       \refstepcounter{\@listctr}\fi
    \quad\@itemlabel\hspace{\labelsep}%
\fi}
\makeatother


\def\f{x+1} \def\g{-x/3+2}  \def\h{-x+3}

\newcommand{\vasymptote}[2][]{
    \draw [densely dashed,#1] ({rel axis cs:0,0} -| {axis cs:#2,0}) -- ({rel axis cs:0,1} -| {axis cs:#2,0});
}

\addpoints
%\printanswers
\noprintanswers

\opengraphsfile{Q3c_Fa18}

\begin{document}
\extrawidth{-0.3in}
\pagestyle{headandfoot}

\setlength{\hoffset}{-.25in}

\extraheadheight{-.3in}
\runningheadrule
\firstpageheader{\bfseries {Precalculus}}{ \bfseries {Quiz 3 }}{\bfseries {9/25/18}} 

\begin{center}
	This quiz has \numquestions\ questions, for a total of \numpoints\
	points and \numbonuspoints\ bonus points.
\end{center}


\firstpagefooter{} %%&&CHANGED
                {}
                {%Points earned: \hbox to 0.5in{\hrulefill}
                % out of  \pointsonpage{\thepage} points
                }
                 
						

\vspace*{0.1cm}
\hbox to \textwidth { \scshape {Name:} \enspace\hrulefill}
\vspace{0.1in}




\pointpoints{point}{points}

\begin{questions}


\addpoints

\question The graph of a function $f(x)$ is shown below.



\begin{center}
	
	\begin{mfpic}[20]{-1}{6}{-1}{5}
		
		%\polyline{(0,-2), (4,1), (4,2), (5,3)}
		
		\polyline{(0,0), (4,4)} 
		
		%\polyline{(2,0), (4,2)}
		
		\point[5pt]{(0,0), (4,4)}
		%\tlabel[cc](4,2.5){\scriptsize $(4,2)$}
		\axes
		
		\xmarks{1,2,3,4,5}
		
		\ymarks{-1,1,2,3,4,}
		
		\tlpointsep{4pt}
		
		\axislabels {x}{{\tiny $1$} 1, {\tiny $2$} 2, {\tiny $3$} 3, {\tiny $4$} 4, {\tiny $5$} 5}
		
		\axislabels {y}{{\tiny $1$} 1,{\tiny $2$} 2, {\tiny $3$} 3, {\tiny $4$} 4,  {\tiny $-1$} -1}
		
		% Grid
		%\drawcolor[gray]{0.005}
		%\gridlines{1, 1}
		\drawcolor[gray]{0.95} 
		\grid{1,1}
		\tlabel[bl](-5,1){Graph of {\bf{f(x)}}}
	\end{mfpic}
	
\end{center}

\begin{parts}
	\part[3]  Fill out the table:
	
	\begin{minipage}{\linewidth}
		\centering
		
		{\setlength{\tabcolsep}{3.1em}  
			{\renewcommand{\arraystretch}{1.2}%
		
		\begin{tabular}{|l|l|l|l|l|}
			\hline
			$x$    & $0$ & $1$ & $2$ & $4$  \\ \hline
			$2f(x)$ &      &     &     &       \\ \hline
				\end{tabular}}} \quad
			\end{minipage}
			\smallskip
			
	%\part[2] Write the domain of $f(2x)$ in interval form. \dotfill
	\part[2] Write the range of $2f(x)$ in interval form. \dotfill 
	\part[2]  The three  transformations listed below are applied (in the given order) to the function $f(x)$. On the grid provided below, draw the graph of the final transformed function $g(x)$ (that is, $g(x)$ is the outcome of applying the three transformations in the given order).

\begin{minipage}[t][][b]{0.4\textwidth}
%\begin{multicols}{2}

		\begin{enumerate}
			\item  Vertical shift down 3 unit, followed by:
			\item Horizontal shift left 2 units, followed by:

			\item  Reflection across the $y$-axis.

			
		\end{enumerate}

		%\end{multicols}
	%(Note: don't forget to draw and label the axes, and include units.)
\end{minipage}
\begin{minipage}[t][][b]{0.6\textwidth}
\begin{center}
	
	\begin{tikzpicture}
	\draw[gray!50, thin, step=0.5] (0,-2) grid (6,2);
	\end{tikzpicture}
	\end{center}
\end{minipage}
	
	

\bonuspart[2] Write the equation of the final transformed function $g(x)$. Show Your work step-by-step.   
\fillwithdottedlines{1.5cm}
\end{parts}
\question[1]  Suppose that $y=t(c)$ represents the function in which $c$ is the number of credits that a student is taking this year $y=t(c)$ is the cost of tuition (in dollars). Suppose a change was made to the tuition so that $y=100+t(c)$ describes the new tuition function. Which of the following situations best describes the new tuition function?

\begin{oneparchoices}
		\choice The tuition has increased by 100\%;
		\choice The total cost of tuition has increased by \$1000;
		\choice The total cost of tuition has increased by \$100;
	\choice The number of credit hours has increased by 100;
	\choice The cost per credit has increased by \$100.

\end{oneparchoices}
\end{questions}

\end{document}                 